\documentclass{ltjsbook}
\usepackage{docmute} 
\usepackage{../../myStyle} 

\begin{document}
\chapter{ベイジアンアプローチと確率的プログラミング1}
それではいよいよ，データ生成モデリングとベイズ推定による実践例を見ていくことにしましょう。
\label{x18_MCMC}
\section{7人の科学者}
ここでは \citet{Lee201709}より「7人の科学者」の例を紹介します。そこでのカバーストーリーは次のようなものです。

\begin{quote}
	実験スキルが大きく異なる7人の科学者が，全員同じ量について測定を行う。彼らが得た数値結果は次の通り。
	$$Y＝{-27.020,3.570,8.191,9.808,9.603,9.945,10.056}$$
	直感的には，最初の二人の科学者はひどく適性を欠いた測定者であり，この量の真の値はおそらく10をわずかに下回るくらいであるように思えるのだが・・・？？
\end{quote}

さあ，これをみて「あれ，どこが統計的な問題なんだ？」と思った人もいるかもしれません。なんらかの測定をして，データのばらつきがあるんだけど，それが測定者によって違うらしい，というのはわかったかと思いますが，これをどうやって統計的な問いにするのでしょうか。まず，ストーリーの背後に「正確な測定値(真の値)はわからないけど，定数のはず」という前提があることを確認しましょう。測定の基本として測定誤差があるので真の値を特定出来はしませんが，測定を繰り返すと真の値に近づいていくはずではあります。また，今回の測定結果は7人それぞれ別の人による測定であることから，測定誤差の出方が測定毎すなわち測定者毎に異なるだろう，という仮説もあります。

これらを踏まえて謎解きをしていきましょう。わからない量を確率の言葉で表現し，データを使って少しでも正解に近づこうとするのが統計のやり方です。
まずここでわからない量は，「正確な測定値」と「個人毎の誤差の大きさ」です。誤差は正規分布に従うと思われますから，データ$Y$は$Y\sim N(\mu, \sigma)$と表すことができます。データ生成メカニズムという意味では，「正規分布に従ってデータが作られる」と言ってもいいかもしれません。また，正確な測定値はここでは$\mu$ということになります\footnote{正確な測定値$\mu$に誤差$N(0,\sigma)$がついてデータになる，すなわち$Y=\mu+N(0,\sigma)$ということですが，誤差の平均は$0$で$\mu$の分だけ正規分布がズレる，そして「分布に従う」を表現する$\sim$を使いたいので，$Y\sim N(\mu,\sigma)$となるのです。}。また，データは7点あって，それぞれ$Y_1,Y_2,\cdots,Y_7$，あるいは一般に$Y_i$と書くことにしましょう。ここで添字の$i$は第$i$回目の測定でもありますし，測定者$i$のことでもあります。この測定者$i$毎に誤差の出方の大きさが変わります。誤差の大小，言い換えれば測定の精度は誤差のSDで表現されていますから，$\sigma$が個人毎に変わる，すなわち$\sigma_i$と書くことができます。これが今回のデータ生成モデルです！

ここまでのことをイメージ図であらわすと図\ref{fig::18_01}のようになります。この図の書き方のポイントは，まずデータを一番下に置くことです。次にこのデータがどこから来たのかな，ということを矢印で表現して書きます。データは誤差を伴うなどで，毎回定数にならないでしょうから，確率分布をその上に書くことになります。確率分布にはその形状を定めるパラメータがあるはずですから，それを書き加えればOKです。


\begin{figure}
	\begin{center}
		\includegraphics[width=10cm,keepaspectratio]{../images/18_MCMC/slide18_01.png}
		\caption{データが生成されるプロセス}
		\label{fig::18_01}
	\end{center}
\end{figure}

しかしここでも，$\mu$と$\sigma_i$は結局なんなんだ，どういう数字なんだということがわからないままです。わからない量は確率の言葉で表現するのがベイジアンの生き様。ここは未知のパラメータがどうやって出てくるかなんてわからないところなので，「わからない」ことを確率で表現する必要があります。パラメータに対して「この辺りにあるだろう」と事前に想定する確率分布のことを\keyterm{事前分布}{じぜんぶんぷ}{prior distribution}というのでした。

ここでは$\mu$は平均0，SD100の正規分布からきていることにしましょう。正規分布というのは左右対称で平均値・中央値・最頻値が一致する単峰の分布です。これは正確な測定値を表しているのでしたから，なんらかの値を取るとしても1つの値だろう，というのは無理のない仮定でしょう。複数の値を取る可能性があるのなら多峰性の分布を仮定したらいいと思いますが，今回はそうではないだろう，と仮定するのです。平均0でSD100ということは，-300から+300の範囲にある可能性が99.7\% だということでもあります\footnote{正規分布は，平均$\pm$1SDの範囲に全体の68\%が，$\pm$2SDの範囲に95\%が，$\pm$3SDの範囲に99.7が入ることが理論的にわかっているのでした。わからない人はRで計算して確認すること！}。実際のデータが-27から10ぐらいの範囲に入っているので，-300から+300の間というのも十分過ぎるぐらい広めの範囲にしてある数字と言えるでしょう。このように，うっすら情報を持っているけどほとんど意味がない程度に縛りをかけるのを，\keyterm{弱情報事前分布}{じゃくじょうほうじぜんぶんぷ}{weakly informative prior distribution}と言います。たとえばこれが身長のデータであれば，$\pm$3mの範囲は十分すぎ，なんなら負の数は取るはずがないので過剰な安全策を置いているとも言えます。それでも「なんらかの事前分布を持つのは，主観的な思い込みだ」というのであれば，$\pm \infty$の一様分布を考えるか，ベイズ法をやめるかということになります。

次に$\sigma_i$です。これは標準偏差パラメータなので負の数を取ることはあり得ません。ここでは\keyterm{半コーシー分布}{はんこーしーぶんぷ}{half-cauchy distribution}を事前分布におくことにします。コーシー分布とは正規分布によくにた形ですが，正規分布より裾が重く\footnote{より極端な値が出る可能性が多いという意味です}，標準偏差(分散)の事前分布に適していると言われている分布です\citep{gelman2006prior}\footnote{この論文では他にも，studentのt分布などを提案していますが，パラメータ数がより少ないコーシー分布を選びました。}。このコーシー分布は正規分布のように左右対称ですが，0を中心に半分に折りたたんで正の値しか取らないように考えたのが半コーシー分布です。今回はスケールパラメータに$5$を置いていますが，特にこの数字に深い意味はありません。気になるようでしたら$10$でも$100$でも変えていただいて結構です\footnote{半コーシー分布がどんな数字になるのか気になる，という人はRで\texttt{rcauchy}関数を使って乱数を発生させ，その挙動を見てみればいいでしょう。乱数発生によって大体の実感を得ることができる，ということを思い出してください。}。

これらをおくことで，モデルの設計図が出来上がりました(図\ref{fig::18_02})。あとはこれをもとに，確率プログラミング言語に書き起こしていけば良いでしょう。
\begin{figure}
	\begin{center}
		\includegraphics[width=12cm,keepaspectratio]{../images/18_MCMC/slide18_02.png}
		\caption{モデルの全体像}
		\label{fig::18_02}
	\end{center}
\end{figure}

\section{Stanコードの書き方}
では図\ref{fig::18_02}という設計図に基づいて，Stanのコードを書いていきましょう。

\subsection{Stanコードの文法}
Stanは事後分布からの乱数を生成するための言語です。その言語仕様はCやC++ににているところがあります。Rで書くのと違うところとして「ブロック単位の記述」「セミコロンによる行の区切り」「変数の宣言」の3つを押さえておきましょう。

\paragraph{ブロック単位の記述}
Stanに書くべきことは\textbf{ブロック}の中に記述します。ブロックは中括弧(\{ ... \})で括ったもので，次の6つの種類があります。
\begin{enumerate}
	\item \texttt{data}ブロック
	\item \texttt{transformed data}ブロック
	\item \texttt{parameters}ブロック
	\item \texttt{transformed parameters}ブロック
	\item \texttt{model}ブロック
	\item \texttt{generated quantities}ブロック
\end{enumerate}

Stanのプログラムに含まれるブロックは，この順番でなければならないと決まっているので注意してください。
この中で，1,3,5番目，すなわち\texttt{data,parameters,model}ブロックが基本的に必要なもので，2.4.6番目のブロックはその派生や応用です。まずは基本ブロックから説明します。

\texttt{data}ブロックはStanと外部とのやりとりをするブロックです。外部からデータを取り込んで，それに基づいて乱数を発生させます。Stanで計算せずに外部から与えるものは全てこのブロックの中で書く必要があります。
\texttt{parameters}ブロックは，パラメータ，すなわちStanで推定したいものを書くところです。Stanは事後分布からの乱数発生機であり，事後分布は確率分布ですからその形状はパラメータによって特徴づけられます。事後分布の形が分からないというのは，事後分布のパラメータがどのような値になっているのか分からない，ということと同じ意味です。その「分からない」「知りたい」パラメータをリストアップしておくブロックです。
\texttt{model}ブロックは，確率モデルすなわち尤度と事前分布を書くところです。ここで書かれた尤度と事前分布が，\texttt{data}ブロックのデータと組み合わさって，事後分布がStan内部で計算され，そこからの乱数が次々取り出されてくることになるのです。

\paragraph{セミコロンによる行の区切り}
このブロックの中に，Stan言語の文を書いていくことになるのですが，ここで1つ覚えておいて欲しいのが，\textbf{Stanの句点はセミコロン}ということです。句点がない文章は読みにくいどころか，機械にとっては意味不明な文字列になりかねませんので，必ず1つの命令・一文がおわったらセミコロンで閉じるようにしてください。これを忘れてエラーになるのが初心者に最も多いミスのひとつです。

\paragraph{変数の宣言}
たとえばRでは数字でも分析結果でも，オブジェクトに代入でき，オブジェクト名は任意につけることができました。またオブジェクトはいつどのタイミングで発生させてもOKで，思いついたらすぐオブジェクトに代入，ということができました。
Stanをはじめ，いくつかの高級言語ではこれに厳格な制限をかけるものがあります。すなわち，どんなオブジェクトを作るかを最初に\keyterm{宣言}{せんげん}{declaration}する必要があるのです。Stanでは使われるオブジェクト(変数)が，どういう数字になるものなのか，どの範囲の数字になるのかをブロックの最初に宣言します。この数字の種類のことを\keyterm{型}{かた}{type}といい，その宣言された数字の型に適合しない数字はエラーとなって弾かれます。なんでそんな窮屈な，と思うかもしれませんが，逆に「ありえない数字は許さない」という意味で，エラーを未然に防いでくれることにもなるのです。

変数の形として，Stanは整数(\keytermL{int型}{intがた})，実数(\keytermL{real型}{realがた})，ベクトル(\keytermL{vector型}{vectorがた})，行列(\keytermL{matrix型}{matrixがた})などがあり，また範囲を\texttt{lower}や\texttt{upper}で指定します。例えば標準偏差は負の数を取らないので，標準偏差の変数を宣言したいときは，\verb|real<lower=0> sig;|のように書きます。
\subsection{Stanコードを書いてみよう}
されこれらを踏まえてコードを書いてみるわけですが，設計図を書いたときにデータを下に，それに紐づく確率分布を上に上に，と下から上に書いていったように，\texttt{model}ブロック$\to$\texttt{parameters}ブロック$\to$\texttt{data}ブロック，と書き進んでいったほうがわかりやすいでしょう。

Stanのコードはメモ帳など，エディタで開くことのできるテキストファイルでいいのですが，拡張子を\texttt{.stan}にしておきましょう。RStudioで書く場合は，\texttt{.stan}にしておくとStanのファイルだと認識してくれるので，ブロックや型宣言の用語などを強調表示してくれるようになります。RStudioでコードを書き始めるときは，ファイル(File)＞新しいファイル(New File)と進むとStanというファイルがあります。これを選ぶとすでにStanのサンプルコードが書かれているファイルが出てきますが，そのファイルの中身は全部消して，次のコード\ref{stan::18_01}のように書き，ファイル名をつけて保存しましょう。ここでは\texttt{sevenScientist.stan} とファイル名をつけたものとします。

\begin{lstlisting}[language=C++,label={stan::18_01},caption={7人の科学者コード}]
data{
  real Y[7];
}

parameters{
  real mu;
  real<lower=0> sig[7];
}

model{
  for(i in 1:7){
    //likelihood
    Y[i] ~ normal(mu,sig[i]);
    //prior
    sig[i] ~ cauchy(0,5);
  }
  //prior
  mu ~ normal(0,100);
}

\end{lstlisting}
\paragraph{コード解説}
\begin{description}
	\item[\texttt{model}ブロック] 尤度と事前分布を書くところ。確率分布からのサンプリングは変数を左に，$\sim$で確率分布につなぎます。\texttt{normal}とか\texttt{cauchy}が確率分布の名前なのでRStudioではハイライト表示されているはず。一行の終わりはセミコロンで。複数のデータや変数，$Y_i$や$\sigma_i$は\verb|Y[i]|や\verb|sig[i]|のように，大括弧(\verb|[ ]| )を後ろにつけて書きます。これは変数が配列(複数の要素を持っている)ときに，その要素を指定する方法です。またデータは7件あって，それぞれについて$i$さんの測定値$Y_i$と測定誤差$\sigma_i$があるので，\texttt{for}文で繰り返しています。変数\verb|sig[i]|についての事前分布も同様に繰り返しています。ちなみに変数名\texttt{mu, sig, Y}などは任意で，予約語でなければ好きな名前にしてもらって構いません。
	\item[\texttt{parameters}ブロック] \texttt{model}ブロックで，変数\texttt{mu}や\texttt{sig}という任意に命名して使っていたわけですが，それは何なのかを宣言してやらねばなりません。これらは求めたいパラメータなのですから，このブロックに型と共に宣言するわけです。変数\verb|mu|は実数型，変数\verb|sig|も実数型ですが7つの要素を持つ配列，しかも負の数を取らないので下限はゼロですよ，という宣言をしているところになります。範囲の設定に\verb|<>|の記号を使っていること注意してください。
	\item[\texttt{data}ブロック] 最後にデータブロックです。モデルの変数でもパラメータでもない，外部から与えら得るものなので，変数名\verb|Y|で宣言しています。実数型で，7人分の要素があります。
\end{description}

いかがでしょうか。書き方(文法)がわかると比較的単純，あるいは率直な書き方をしていることがお分かりいただけるかと思います。尤度と事前分布とデータを書くだけで，Stanは事後分布からの乱数を生成してくれます。確率密度関数を書いたり解いたりしなくても，結果の近似計算ができるというのはとても便利なことではないでしょうか。
\section{Stanを使ったMCMCの実践}
さてStanファイルが準備できましたから，これは一旦保存しておいて，今度はこれをRの方から利用してやることになります。
Rからは\texttt{rstan}パッケージあるいは\texttt{cmdstanr}パッケージを使って，Stanを呼び出して結果を返してもらう，という流れです。Stanの言語は一度機械語に翻訳(コンパイル)されるので，その間はしばらく待つ必要がありますが，それが終わるとStanから事後分布の乱数が次々生成され，Rに返ってきます。Rはそれら乱数からのサンプルを事後分布の代表値からなるデータセットとして扱って，記述統計や可視化を通じて確率分布の解釈に使うのですね。

それではどのようにRから呼び出すのか，実際のコードを見てみましょう。呼び出し方は\texttt{rstan}パッケージと\texttt{cmdstanr}パッケージのどちらを使うかによって少し異なります。共通部分のあと，両方の呼び出し方について記述しますので，自分の環境にあった方を試してください。
\paragraph{共通部分}
どのパッケージを使うにしても，まずは次の5行分を書いておいてください(code\ref{code::18_02})。

\begin{lstlisting}[language=C++,label={code::18_02},caption={環境の準備と共通部分}]
rm(list = ls())
library(tidyverse)
library(bayesplot)
# データ
x <- c(-27.020, 3.570, 8.191, 9.898, 9.603, 9.945, 10.056)
\end{lstlisting}

\paragraph{rstanパッケージによる呼び出し}

\texttt{rstan}パッケージを使う場合は次のように書きます(code\ref{code::18_03})。

\begin{lstlisting}[language=C++,label={code::18_03},caption={rstanパッケージによる呼び出し}]
library(rstan)
options(mc.cores = parallel::detectCores())
rstan_options(auto_write = TRUE)
model <- stan_model("sevenScientist.stan")
fit1rstan <- sampling(model, data = list(Y = x))
sample <- fit1rstan %>%
  rstan::extract() %>%
  as.data.frame() %>%
  as_tibble()
\end{lstlisting}
\paragraph{コード解説}
\begin{description}
	\item[1行目]　パッケージの呼び出し
	\item[2-3行目] CPUの並列化や上書き保存など，パッケージの設定
	\item[4行目] モデルのコンパイルを実行して\texttt{model}オブジェクトに入れます。コンパイルには少し時間がかかります。
	\item[5行目] 乱数発生。ここでRのデータをStanに渡しています。実行すると画面に色々表示されるかと思います。
	\item[6-9行目] 得られた乱数をデータセットに整形。まず\texttt{rstan::extract}でMCMCサンプルだけを取り出し，\texttt{as.data.frame}でデータフレーム型に，\verb|as_tibble|でtibble型に変換します。
\end{description}

\paragraph{cmdstanrパッケージによる呼び出し}

\texttt{cmdstanr}パッケージを使う場合は次のように書きます(code\ref{code::18_04})。

\begin{lstlisting}[language=C++,label={code::18_04},caption={cmdstanrパッケージによる呼び出し}]
library(cmdstanr)
model <- cmdstan_model("sevenScientist.stan")
fit1cmdstan <- model$sample(
  data = list(Y = x),
  chains = 4,
  parallel_chains = 4
)
sample <- fit1cmdstan$draws() %>% 
  posterior::as_draws_df() %>% 
  as_tibble()
\end{lstlisting}
\paragraph{コード解説}
\begin{description}
	\item[1行目] パッケージの呼び出し
	\item[2行目] モデルのコンパイルを実行して\texttt{model}オブジェクトに入れます。コンパイルには少し時間がかかります。
	\item[4-8行目] 乱数発生。ここでRのデータをStanに渡したり，並列化の指定をしたりします。実行すると画面に色々表示されるかと思います。
	\item[9-11行目] 得られた乱数をデータセットに整形。まず\verb|posterior::as_draws_df|でMCMCサンプルだけを取り出し，\verb|as_tibble|でtibble型に変換します。
\end{description}

どちらのパッケージで実行してもらっても結構ですが，コンパイル$\to$サンプリングという流れは同じです。またサンプリングの際に，Rからデータを渡すことになりますが，データはR環境では\texttt{x}というオブジェクト名だったものを，Stanの\texttt{data}ブロックで宣言した\texttt{Y}という変数にして渡すのだ，ということが書かれています。データは複数のものを与えることがあるので，\texttt{list}関数で並べて渡すことにします。そのほかの設定については，順次説明します。

最終的に\verb|tibble|型でえられたMCMCサンプルのデータセットが出来上がります。これはパッケージによって少し変数名が異なりますが，内容は同じで，今回の設定だと4000行のデータになっているはずです。つまり事後分布からの乱数を4000個発生させたことになります。

そしてこの各行が1回のサンプリングで得られた\textbf{事後分布からの代表値}です。一部を見てみましょう(出力\ref{RS:out18_01})。
1回目のサンプリングで，$9.98,23.4,4.63,1.13,0.123,5.17,0.0264,1.24,-9.10$という数字のセットがあります。変数名にあるように，前から$\mu, \sigma_1,\sigma_2,\cdots,\sigma_7$の代表値と，\verb|lp__|の値が得られています。最後の\verb|lp__|は\keyterm{対数尤度}{たいすうゆうど}{log posterior}の略で，ここでは正確には事後対数尤度と呼ばれるものです。今回の代表値から計算された尤度の対数ということで，パラメータの推定値とは違いますのでここでは気にしないことにしましょう。ということで改めて，1回目のサンプリングで$9.98,23.4,4.63,1.13,0.123,5.17,0.0264,1.24$という数字のセットが得られました。このように，今回求めたかったパラメータは8つあったわけですが，その8次元からなる空間の1つの座標点を取り出した1回目，ということです。2回目は$9.90,27.2,13.1,2.86,1.23,0.779,0.0449,0.446$というセットでした。以下3,4,..と続き，4000サンプルを得たわけです。このように結果は8次元からなるパラメータの\keyterm{同時確率空間}{どうじかくりつくうかん}{joint probability space}から得られており，一行一行が1回1回のサンプリングステップなのです。

\begin{Rscreen}{MCMCサンプルの結果(ごく一部)}{out18_01}
	\begin{verbatim}
        # A tibble: 4,000 × 9
        mu sig.1 sig.2 sig.3 sig.4 sig.5   sig.6 sig.7   lp__
     <dbl> <dbl> <dbl> <dbl> <dbl> <dbl>   <dbl> <dbl>  <dbl>
   1  9.98  23.4  4.63  1.13 0.123 5.17   0.0264 1.24   -9.10
   2  9.90  27.2 13.1   2.86 1.23  0.779  0.0449 0.446  -7.72
   3  9.99  20.9  5.12  1.11 0.132 4.48   0.0397 0.937  -8.97
\end{verbatim}
\end{Rscreen}

\section{MCMC結果の診断}
さて，MCMCが実行できましたが，その中身を少しみてみましょう。サンプリング結果を代入したオブジェクト名をそのまま入力すると，要約された結果が示されます。パッケージによって出力例は少し異なりますが，ほぼ同じ情報が含まれています。

\begin{SCrstan}{MCMCの結果出力(rstanパッケージの場合)}{out18_02}
	\begin{verbatim}
Inference for Stan model: sevenScientist.
4 chains, each with iter=2000; warmup=1000; thin=1; 
post-warmup draws per chain=1000, total post-warmup draws=4000.

        mean se_mean    sd   2.5%    25%   50%   75%  97.5% n_eff Rhat
mu      9.74    0.02  0.60   8.11   9.62  9.89  9.98  10.52   607 1.02
sig[1] 38.61    1.37 43.68  14.36  23.33 27.16 41.93 123.28  1016 1.01
sig[2]  9.44    0.57 12.88   2.86   4.63  6.22 10.57  33.54   515 1.00
sig[3]  4.64    0.40 14.66   0.63   1.33  2.49  4.98  17.98  1354 1.01
sig[4]  1.96    0.18  4.14   0.04   0.17  0.79  2.18  10.41   511 1.01
sig[5]  3.35    0.23  7.54   0.11   0.73  1.87  4.56  13.23  1044 1.01
sig[6]  2.00    0.25  3.65   0.03   0.22  0.77  2.44  10.55   216 1.02
sig[7]  2.50    0.17  5.71   0.09   0.44  1.24  2.49  12.28  1133 1.01
lp__   -9.45    0.07  1.67 -13.80 -10.19 -9.10 -8.39  -7.07   570 1.01
    \end{verbatim}
\end{SCrstan}

\begin{SCcmd}{MCMCの結果出力(cmdstanrパッケージの場合)}{out18_03}
	\begin{verbatim}
 variable  mean median     sd   mad     q5    q95 rhat ess_bulk ess_tail
   lp__   -9.70  -9.35   1.91  1.69 -13.28  -7.30 1.00      914     1694
   mu      9.64   9.81   0.67  0.36   8.28  10.36 1.01      879     1101
   sig[1] 54.19  31.91 105.83 16.95  15.41 142.60 1.01      663      215
   sig[2] 10.48   6.71  20.30  3.96   2.88  25.82 1.00     1272      935
   sig[3]  4.30   2.76   6.19  2.01   0.67  12.16 1.01     1281     1104
   sig[4]  2.52   1.31   3.88  1.47   0.11   8.58 1.01      603      252
   sig[5]  2.68   1.37   5.65  1.39   0.19   9.02 1.00     1291     1535
   sig[6]  2.66   1.26   5.11  1.47   0.09   9.46 1.01      492      312
   sig[7]  2.87   1.42   5.64  1.56   0.15   9.49 1.01      904     1049
    \end{verbatim}
\end{SCcmd}

ここで確認して欲しいのは，\keytermL{Rhat}{Rはっと}です。これは「複数チェインの収束の程度」を表しています。
まずMCMCによってサンプルを得る方法は乱数だったことを思い出してください。乱数ですから毎回同じ数字からスタートするわけではありません。
そしてある乱数は次の数字を生むステップになっており，一歩一歩，事後分布の形に近づいていくものです。しかし違う数字からスタートしても，目標とする事後分布は同じはずですから，最終的にはこのステップの鎖はよく絡まったもの，つまりどこからスタートしても同じあたりをウロウロするものになるはずです。ですから，MCMCがきちんとできているかどうか，事後分布を正しく代表した数字になっているかどうかを確認する必要があります。この\keytermL{Rhat}{Rはっと}という数字は，$1.0$であればピッタリ一致という数字で，$1.1$よりも小さければ十分「絡まっている」と判断できます。今回はいずれも十分な収束が得られたと言えるでしょう。

このことを視覚的に確認することもできます。\keyterm{トレースプロット}{トレースプロット}{trace plot}というのがそれで，MCMCのステップを可視化し，十分に絡まっていることを見ておくのです(図\ref{fig::18_03})。
ちなみにこの図，\texttt{rstan}パッケージの場合は\texttt{traceplot(fit1rstan)}，\texttt{cmdstanr}の場合は\verb|fit1cmdstan$draws("mu") %>% bayesplot::mcmc_trace()|とすれば描画できます。

\begin{figure}
	\begin{center}
		\includegraphics[width=10cm,keepaspectratio]{../images/18_MCMC/Rplot18_01.png}
		\caption{トレースプロットの例}
		\label{fig::18_03}
	\end{center}
\end{figure}
これについては悪い例を見た方が早いかもしれません。図\ref{fig::18_04}に，チェインが絡まない場合のトレースプロットを示しました。このように，複数のチェインがそれぞれバラバラ，迷子になっている場合は，手元に生成された乱数が事後分布の適切な代表になっていないと判断することになります。
\begin{figure}
	\begin{center}
		\includegraphics[width=10cm,keepaspectratio]{../images/18_MCMC/Rplot18_02.png}
		\caption{チェインが絡まない場合}
		\label{fig::18_04}
	\end{center}
\end{figure}

次に，\keyterm{有効サンプルサイズ}{ゆうこうさんぷるさいず}{effective sample size}を見ておきましょう。\texttt{rstan}パッケージを使っている例の場合は\verb|n_eff|，\texttt{cmdstanr}パッケージを使っている例の場合は\verb|ess_tail|を見ましょう。これはMCMCサンプルがうまく取れたかどうかを診断した結果で，ここが少なくとも3桁ほどなければ不十分なのだ，と考えます。いくらあったらいいかについては，どれだけMCMCサンプルを発生させたかに依存するので一概には言えませんが，少なくとも効果的に得られているのが100サンプルぐらいないとダメなんだ，ということです。

ちなみに，MCMCサンプルのチェイン(発生させた系列数)や，MCMCサンプルの数は，関数のオプションで指定します。今回は指定なしだったので，デフォルトの値で4チェイン，各々1000サンプルずつ作るようになっていました。これらを明示的に設定するには次のようにします(code\ref{code::18_05},code\ref{code::18_06},)。
\begin{lstlisting}[language=C++,label={code::18_05},caption={rstanパッケージによるオプション指定}]
fit1rstan <- sampling(model,
  data = list(Y = x),
  chains = 4, iter = 5000, warmup = 1000
)
\end{lstlisting}

\begin{lstlisting}[language=C++,label={code::18_06},caption={cmdstanrパッケージによるオプション指定}]
fit1cmdstan <- model$sample(
  data = list(Y = x),
  chains = 4,
  parallel_chains = 4,
  iter_warmup = 1000,
  iter_sampling = 4000
)
\end{lstlisting}

ここにある\texttt{chains}というのが生成するMCMCサンプルの連鎖系列の数です。「どこからスタートしても同じ事後分布にたどり着く」ことを検証するために，少なくとも複数のサンプルが必要で，デフォルトでは4本のチェインを作っていることになります。\texttt{cmdstanr}パッケージの方は，並列化する本数も別途\verb|parallel_chains|関数で指定します。

また生成する乱数の数ですが，ここに\keyterm{ウォームアップ}{うぉーむあっぷ}{wamup}という単語が出てきています。これはMCMCサンプルの際，最初の数ステップは調整のための準備運動期間中なので，この期間に得られたMCMCサンプルは代表値として使わない，という意味です。\texttt{rstan}パッケージの場合，全体の反復(iter)回数からウォームアップの分を引き算して，残りの期間をチェイン数発生させることになります。今回の例では$(5000-1000)\times 4=16000$サンプル出来上がることになります。

\texttt{cmdstanr}パッケージの場合は引き算ではなく，最初の1000回は捨てて，サンプリング期間$4000\times 4=16000$で計算できます。

\section{MCMCの結果の解釈}
さあ長々と話をしてきましたが，どうやらRhatや有効サンプルサイズにも問題がなかったので，MCMCは事後分布からの適切な代表値を産んでいたと考えられますから，その結果を考えてみましょうか。

MCMCサンプルを取り出して作ったデータセットから，code\ref{code::18_07}のようなコードを書いてみましょう(前回作成した\verb|map_estimation|関数を読み込んでおく必要があります)。
\texttt{rstan}パッケージの結果は各行が1回のサンプリングに該当しますが，それを明確にするためにまず行番号を\texttt{iter}変数として生成し，そのあとで並べ替えなどを行なっています。

\begin{lstlisting}[language=C++,label={code::18_07},caption={rstanパッケージによる結果の要約}]
sample_rstan %>% 
    rowid_to_column("iter") %>%
    dplyr::select(-lp__) %>% 
    pivot_longer(-iter) %>% 
    group_by(name) %>% 
    summarise(
        EAP = mean(value),
        MED = median(value),
        MAP = map_estimation(value),
        U95 = quantile(value,prob=0.975),
        L95 = quantile(value,prob=0.025)
    )
\end{lstlisting}

\texttt{cmdstanr}パッケージの結果はチェイン番号，各チェインでの反復回数，延反復回数が変数として記録されているので，剰余変数を除外するだけでOKです(code\ref{code::18_08})。

\begin{lstlisting}[language=C++,label={code::18_08},caption={cmdstanrパッケージによる結果の要約}]
sample_cmdstan %>% 
    dplyr::select(-lp__,-.chain,-.iteration) %>% 
    pivot_longer(-.draw) %>% 
    group_by(name) %>% 
    summarise(
        EAP = mean(value),
        MED = median(value),
        MAP = map_estimation(value),
        U95 = quantile(value,prob=0.975),
        L95 = quantile(value,prob=0.025)
    )
\end{lstlisting}

もうこのコードの意味の説明は省略していいですね。この結果示される複数の代表値がありますが，私の結果では表\ref{tbl::18_01}のようになりました。


\begin{table}[ht]
	\caption{MCMCの結果}
	\centering
	\begin{tabular}{l|ccccc}
		\hline
		パラメータ & EAP    & MED    & MAP    & U95     & L95    \\
		\hline
		$\mu$      & 9.570  & 9.622  & 9.596  & 10.547  & 7.753  \\
		$\sigma_1$ & 42.197 & 26.955 & 23.101 & 151.573 & 14.259 \\
		$\sigma_2$ & 8.776  & 6.853  & 6.849  & 30.356  & 2.590  \\
		$\sigma_3$ & 4.310  & 3.232  & 3.626  & 17.280  & 0.440  \\
		$\sigma_4$ & 2.756  & 1.102  & 0.993  & 15.750  & 0.085  \\
		$\sigma_5$ & 2.507  & 1.052  & 0.358  & 11.823  & 0.022  \\
		$\sigma_6$ & 2.503  & 1.676  & 2.021  & 12.339  & 0.076  \\
		$\sigma_7$ & 2.598  & 1.020  & 0.674  & 13.548  & 0.100  \\\hline
	\end{tabular}
	\label{tbl::18_01}
\end{table}

ここでは\keyterm{事後確率最大値}{じごかくりつさいだいち}{Maximum A Posterior}を見ることにしましょう。
これによると，$\mu$すなわち真の測定値は9.596ぐらいじゃないか，という推測がなされています。10より少し小さいくらいではないか，と思っていたようですが，こうして計算してみると確かにそれぐらいのようですね。
そして7人それぞれの測定精度が$\sigma_i$として算出されています。これをみると，最初の測定者は23.101，2番目の人は6.849になっています。これが大きいということは，誤差の幅広さを表しているのですから，測定精度が悪いということを意味します。確かに最初の二人は非常に精度が悪く，他のベテラン測定者が0.3とか0.9ぐらいの誤差であることに比べれば，圧倒的に悪いということが言えますね。

もちろんこれは点推定値で，真の測定値が9.596に違いないとか，最初の人の測定精度が23.101だと断言すると，おそらくほぼ確実に外れた予測ということになるでしょう。そこで幅を持った予測，すなわち\keyterm{区間推定}{くかんすいてい}{interval estimation}を行えば良いのです。区間を考えれば，真の測定値は10.547から7.753の範囲にあるだろうとか，最初の測定者の精度も最悪151.573，最善なら14.259ぐらいである可能性があるわけです。

ベイズ推定のポイントとして，これらは事後分布，すなわち\textbf{確率分布による予測}であるということが挙げられます。真の測定値$\mu$のありそうな範囲が10.547から7.753にある確率が95\%だ，という言い方ができるところです。モーメント法による\keyterm{信頼区間}{しんらいくかん}{confidential intervals}とは違い，ベイズ法の予測は\keyterm{確信区間}{かくしんくかん}{Credible Intervals}という言い方になります\footnote{信頼区間の場合は，ここにあると予測したとして，その予測が当たる確率を表しています。すなわち当たるか外れるか，0/1の話をしているのであって，この「範囲に存在する」という大きさ・幅の話をしているわけではないことに注意です。}。

いかがでしたでしょうか。今回はモデリングとベイズ推測の実際，Stanコードの書き方を実践してみました。
この「7人の科学者」が面白い点は，少ないサンプル(たった7件のデータ！)で推測ができること，簡単なコードで検証できること，というのもありますが，なにより標準偏差(分散)を検証対象とすることにあると私は思います。心理学の実験のほとんどが，平均の比較，操作の効果の話になっているのですが，データの散らばりというのも測定精度のように心理学的に意味のあるものとして考えられ，検証の対象にしたっていいのです。もちろん正規分布でなくてもいいですし，様々な分布の様々なパラメータに意味があるなら，それをわからないものとして検証しようと言える，というのはとても可能性が広がる話ではないでしょうか。
\section{課題}
次の計算をするR/Stanコードを記述し，提出してください。なお提出されたコード単体でバグがなく動くことが確認できないものは，未提出扱いになります。コードの書き方などわからないところがあれば，曜日別TAか小杉までメールで連絡し，指導を受けてください。
\begin{enumerate}
	\item 8人目の科学者が現れて，測定値18.25だったと報告してきました。この人のデータも含めた分析モデルにするために，RおよびStanのコードを書き換え，MCMC推定してください。また7人の時とくらべて結果がどう変わったでしょうか？気づいたことを報告してください。
	\item MCMCサンプルの精度を上げるため，5つのチェインで，warmup期間を10000，最終的なMCMCサンプルを10万点得られるようにRのコードを書き換えてください。
	\item MCMCサンプルを使って，$\mu$と$\sigma$の分布の形を可視化するコードを書いてください。
\end{enumerate}

\end{document}
